\makeatletter
\def\input@path{{../}{../../}{../../../}} % add as many parents as you need
\makeatother
\input{pset_template.tex}

\begin{document}

\section*{Lecture 4 Summary: Policy Gradient Methods}

\subsection*{1. Motivation}

Value-based methods (e.g., Value Iteration, Q-Learning) learn value functions and then act greedily.
However, they can become unstable or inefficient with nonlinear function approximation.
Policy-gradient methods instead \emph{directly} parameterize and optimize the policy itself:
\[
		J(\theta) = \mathbb{E}_{\pi_\theta}\!\left[ \sum_{t=0}^{\infty} \gamma^t r_t \right].
\]
The goal is to find parameters $\theta$ that maximize $J(\theta)$.

\subsection*{2. Stochastic Policy Parameterizations}

\begin{table}[h!]
		\centering
		\begin{tabular}{lll}
				\toprule
				\textbf{Setting} & \textbf{Example} & \textbf{Log-gradient} \\
				\midrule
				\textbf{Discrete actions} &
				Softmax over features $f_\theta(s,a)$:
				\[
						\pi_\theta(a|s) = \frac{e^{f_\theta(s,a)}}{\sum_{a'} e^{f_\theta(s,a')}}
				\]
				&
				\[
						\nabla_\theta \log \pi_\theta(a|s)
						= \nabla_\theta f_\theta(s,a)
						- \sum_{a'} \pi_\theta(a'|s)\,\nabla_\theta f_\theta(s,a')
				\]
				\\[1.2em]
				\textbf{Continuous actions} &
				Gaussian policy $\mathcal{N}(\mu_\theta(s), \sigma^2)$ &
				\[
						\nabla_\theta \log \pi_\theta(a|s)
						= \frac{(a - \mu_\theta(s))}{\sigma^2}\,\nabla_\theta \mu_\theta(s)
				\]
				\\
				\bottomrule
		\end{tabular}
		\caption{Common stochastic policy parameterizations and their log-gradients.}
\end{table}

These forms make $\nabla_\theta \log \pi_\theta(a|s)$ directly computable by autodiff.

\subsection*{3. Policy Gradient Theorem}

The \emph{Policy Gradient Theorem} gives a simple, model-free expression for the gradient:
\[
		\nabla_\theta J(\theta)
		= \mathbb{E}_{s \sim d^{\pi_\theta},\, a \sim \pi_\theta}
		\!\big[ Q^{\pi_\theta}(s,a)\,\nabla_\theta \log \pi_\theta(a|s) \big].
\]
Key insight: the derivative does not include $\nabla_\theta d^{\pi_\theta}(s)$,
so we can estimate this gradient using on-policy samples.

\subsection*{4. REINFORCE Algorithm (Williams, 1992)}

A Monte Carlo estimator of the gradient:
\[
		\hat{g} = \sum_{t=0}^{T-1} R_t \, \nabla_\theta \log \pi_\theta(a_t|s_t),
		\quad
		R_t = \sum_{t' \ge t} \gamma^{t'-t} r_{t'}.
\]
\begin{algorithm}[H]
		\caption{REINFORCE}
		\begin{algorithmic}
				\State Initialize policy parameters $\theta$
				\For{each episode}
				\State Sample trajectory $\tau = (s_0,a_0,r_0,\ldots,s_T)$ under $\pi_\theta$
				\For{each timestep $t$ in $\tau$}
    \State Compute return $R_t = \sum_{t' \ge t} \gamma^{t'-t} r_{t'}$
    \State $\theta \leftarrow \theta + \alpha R_t \nabla_\theta \log \pi_\theta(a_t|s_t)$
				\EndFor
				\EndFor
		\end{algorithmic}
\end{algorithm}

REINFORCE is unbiased but has high variance.

\subsection*{5. Variance Reduction via Baselines}

Subtracting a baseline $b(s)$ that does not depend on $a$ leaves the gradient unbiased:
\[
		\nabla_\theta J(\theta)
		= \mathbb{E}\!\left[ (Q^\pi(s,a) - b(s)) \, \nabla_\theta \log \pi_\theta(a|s) \right].
\]
Choosing $b(s) = V^\pi(s)$ minimizes variance and defines the \emph{Advantage}:
\[
		A^\pi(s,a) = Q^\pi(s,a) - V^\pi(s).
\]
The update then becomes:
\[
		\theta \leftarrow \theta + \alpha A_t \, \nabla_\theta \log \pi_\theta(a_t|s_t).
\]

\subsection*{6. Actor--Critic Methods}

In practice, $Q^\pi$ or $V^\pi$ are unknown and learned with a \emph{critic}.
The critic minimizes TD-error:
\[
		\delta_t = r_t + \gamma V_\omega(s_{t+1}) - V_\omega(s_t),
		\qquad
		\omega \leftarrow \omega - \beta \, \delta_t \nabla_\omega V_\omega(s_t).
\]
The actor then updates using $\delta_t$ as a low-variance estimate of the advantage:
\[
		\theta \leftarrow \theta + \alpha \, \delta_t \, \nabla_\theta \log \pi_\theta(a_t|s_t).
\]

\subsection*{7. Average Reward (Undiscounted) Case}

For continuing tasks, the objective is the average steady-state reward:
\[
		\rho^\pi = \sum_s d^\pi(s) \sum_a \pi(a|s) R(s,a),
\]
and the gradient takes the same form:
\[
		\nabla_\theta \rho^\pi
		= \mathbb{E}_{s,a \sim \pi}\!\left[
				A^\pi(s,a)\,\nabla_\theta \log \pi_\theta(a|s)
		\right].
\]

\subsection*{8. Conservative and Trust-Region Policy Updates}

Large policy changes can degrade performance.
To guarantee monotonic improvement:
\[
		\pi_{k+1} = (1-\alpha)\pi_k + \alpha \pi^*_k,
\]
or impose a KL constraint (as in TRPO):
\[
		\max_\theta
		\mathbb{E}\left[
				\frac{\pi_\theta(a|s)}{\pi_{\theta_{\text{old}}}(a|s)} A^{\pi_{\theta_{\text{old}}}}(s,a)
		\right]
		\text{ subject to }
		\mathbb{E}\!\left[ \mathrm{KL}(\pi_{\theta_{\text{old}}}\,\|\,\pi_\theta) \right] \le \delta.
\]
The practical approximation is PPO’s \emph{clipped surrogate} objective.

\subsection*{9. Key Takeaways}

\begin{itemize}
		\item Policy Gradient Theorem enables model-free differentiation of expected return.
		\item REINFORCE is unbiased but high-variance; baselines reduce variance.
		\item Actor--Critic uses bootstrapped value estimates to stabilize learning.
		\item TRPO/PPO constrain policy updates for reliable improvement.
		\item Policy-based methods handle continuous actions and stochasticity naturally.
\end{itemize}

\noindent
\textbf{In one line:}
Policy gradient methods learn by differentiating through the policy itself;
REINFORCE provides the estimator, baselines and critics control variance,
and trust-region methods make deep RL stable.

\end{document}
